% !TeX root = ../main.tex
\section{模型建立}

除特别说明外，用户集合、请求集合及可用路径网络均指当前优化轮对应的集合，其轮次索引予以省略。各集合随滚动过程更新，用户与请求的标识保持不变。

将运营时间划分为长度为 $\delta$ 的时间段，记
\begin{subequations}\label{eq:time_grid}
\begin{align}
t_n&=n\delta,\label{eq:time_points}\\
\mathcal I_n&=[t_n,t_{n+1}),\label{eq:time_intervals}\\
\mathcal T_\ell&=\{\ell,\ldots,\ell+H-1\}.\label{eq:prediction_indices}
\end{align}
\end{subequations}
其中 $\ell$ 为当前轮，$H$ 为预测时间段数。本轮预测范围为 $[t_\ell,t_{\ell+H})$，在 $t_n$ 处理服务及换电，随后在 $\mathcal I_n$ 内充电。每轮只执行当前时间段，并在下一时刻使用新观测重新优化。默认 $\delta=5$~min，能量计算时取 $\delta=1/12$~h。

充电功率每段更新，换电路径每 $K$ 段允许更新一次。以运营起点为共同基准，定义已知参数
\begin{subequations}\label{eq:update_clock}
\begin{align}
g_\ell&=\ind\{\ell\bmod K=0\},\label{eq:path_update_indicator}\\
K&=3.\label{eq:path_update_interval}
\end{align}
\end{subequations}
因此，两类预约用户的路径均每 $K\delta=15$~min 允许更新一次。$g_\ell=1$ 时联合优化全部预约用户的路径与充电功率；$g_\ell=0$ 时，在途用户保持最近发布的换电安排，尚未进入高速公路的用户保持上一次优化后保留的计划，仅优化充电功率。未进入高速的用户首次参与优化时，以日前路径作为初始计划。进入高速前，优化后的计划不向用户发布，也不产生实际调整费用。图~\ref{fig:updates} 示意两种决策的更新安排。

\begin{figure}[!htbp]
\centering
\begin{tikzpicture}[x=2.55cm,y=0.7cm,font=\small]
\draw[-{Latex[length=2mm]},thick] (0,0)--(5.15,0);
\foreach \x/\lab in {0/0,1/5,2/10,3/15,4/20,5/25}{
 \draw (\x,-0.08)--(\x,0.08);
 \node[above] at (\x,0.13) {\lab~min};
 \fill[teal!70!black] (\x,-0.7) circle (2.3pt);
}
\foreach \x in {0,3}{\fill[orange!90!black] (\x,-1.5) rectangle ++(0.10,0.16);}
\node[anchor=east] at (-0.14,-0.7) {充电};
\node[anchor=east] at (-0.14,-1.4) {路径};
\node[below,align=center] at (2.5,-1.9) {每轮更新车辆、电池和请求状态，并执行当前时间段决策};
\end{tikzpicture}
\caption{充电功率与换电路径的更新安排}
\label{fig:updates}
\end{figure}

表~\ref{tab:notation} 按集合与索引、输入参数、决策变量、由决策计算的量与价值函数列出主要符号。预约用户按 O--D 对分组，以 $\mathcal P$ 表示 O--D 对集合，$p\in\mathcal P$ 为组索引，$k$ 为该组内的用户索引。用户由 $(p,k)$ 唯一标识，编号在滚动过程中保持不变。到站请求统一以 $r$ 标识；某一请求所在站点记为 $i(r)$。费用和价格参数统一采用小写 $c$，合计费用采用大写 $C$，服务收入记为 $I$。

\begin{longtable}{@{}p{4.4cm}p{11.1cm}@{}}
\caption{主要符号}\label{tab:notation}\\
\toprule 符号 & 含义 \\ \midrule
\endfirsthead
\caption[]{主要符号（续）}\\
\toprule 符号 & 含义 \\ \midrule
\endhead
\bottomrule\endfoot
\multicolumn{2}{@{}l@{}}{\textbf{集合与索引}}\\*
$\mathcal S;\ i,j$ & 换电站集合及站点索引 \\
$\mathcal B_i;\ b$ & 站点 $i$ 的电池充电槽集合及充电槽索引 \\
$\mathcal P;\ p$ & O--D 对集合及索引 \\
$\mathcal K^{\mathrm{fut},p},\mathcal K^{\mathrm{enr},p}$ & O--D 对 $p$ 下本轮尚未进入、已经进入高速公路的预约用户集合 \\
$\mathcal K^p;\ k$ & O--D 对 $p$ 下本轮预约用户集合及组内用户索引, $\mathcal K^p=\mathcal K^{\mathrm{fut},p}\cup\mathcal K^{\mathrm{enr},p}$ \\
$\R,\R_i;\ r$ & 本轮全部请求集合、站点 $i$ 的请求集合及请求索引 \\
$\R_i^A,\R_i^R$ & 站点 $i$ 的预约请求集合与随机请求集合 \\
$\mathcal V^{p,k},\mathcal A^{p,k}$ & 用户 $(p,k)$ 本轮路径网络的节点集合与可用弧集 \\
$\mathcal T_\ell;\ n,\ell$ & 本轮预测时间段集合；$n$ 为时段索引，$\ell$ 为当前轮索引 \\
\midrule
\multicolumn{2}{@{}l@{}}{\textbf{输入参数}}\\*
$\delta$ & 时间段长度 \\
$H,K$ & 预测域长度与路径更新间隔，均以时间段数计 \\
$o^{p,k},d^p$ & 用户 $(p,k)$ 本轮路径的起点与出口节点 \\
$v_{i,j}^p$ & O--D 对 $p$ 中节点 $i$ 至 $j$ 的行驶 SOC 消耗 \\
$E_B$ & 电池容量（kWh） \\
$\eta_i$ & 站点 $i$ 的充电效率 \\
$\rho_r$ & 请求 $r$ 对应的退回电池 SOC \\
$S_{i,b}^{\mathrm{obs}}$ & 本轮开始时充电槽 $b$ 中电池的观测 SOC \\
$\overline P_{i,b}^{\mathrm{ch}},\overline P_i^{\mathrm{sta}}$ & 设备充电功率上限与站点总充电功率上限（kW） \\
$\bar t_r^{\mathrm{arr}}$ & 当前等待请求的实际到站时刻，或不需要先在本轮完成另一次换电的请求的预计到站时刻 \\
$\tau_{j,i}^{p,k}$ & 用户 $(p,k)$ 从换电站 $j$ 到下一换电站 $i$ 的预测行驶时间，本轮求解时已知且为正 \\
$\Delta$ & 请求自到站起允许的最大等待时间 \\
$\bar y_{i,j}^{p,k}$ & 用户 $(p,k)$ 的日前换电路径是否使用弧 $(i,j)$ \\
$c_{i,n}^{\mathrm{ch}},c_{i,n}^{\mathrm{sw}}$ & 充电电价与换电服务单价（元/kWh） \\
$c^{\mathrm{adj}},c^{\mathrm{fail}}$ & 单次路径调整惩罚与每个未满足预约请求的惩罚（元/次） \\
$\bar x_i^{p,k},\widehat x_i^{p,k},x_i^{p,k,\mathrm{pub}}$ & 本轮开始时已知的日前路径、保留计划和最近发布路径对应的换电站安排 \\
$\bar{\boldsymbol c}_r,\boldsymbol c_{r,m}$ & 本轮求解前根据车辆行驶预测计算的候选终端车辆信息，见第~\ref{sec:value}~节 \\
$\beta$ & 终端价值权重\\
\midrule
\multicolumn{2}{@{}l@{}}{\textbf{决策变量}}\\*
$y_{i,j}^{p,k}$ & 用户 $(p,k)$ 的本轮换电路径是否经过弧 $(i,j)$，二元变量 \\
$P_{i,b,n}$ & 站点 $i$ 的设备 $b$ 在时间段 $n$ 的充电功率（kW） \\
$\alpha_{b,r,n}$ & 是否在 $t_n$ 用充电槽 $b$ 的电池服务请求 $r$，二元变量 \\
\midrule
\multicolumn{2}{@{}l@{}}{\textbf{由决策计算的量与价值函数}}\\*
$S_{i,b,n}^{-}$ & 充电槽 $b$ 中电池在 $t_n$ 当次换电前的 SOC，由初始电量、换电安排和充电功率确定 \\
$\chi_{p,k}$ & 本轮的换电站安排是否与日前路径（未进入高速）或最近发布路径（在途）不同，由路径比较确定 \\
$f_r$ & 预约请求 $r$ 是否在本轮预测期内超过等待期限仍未获服务，由是否需要服务、截止时刻和换电次数确定 \\
$Q_{r,n}$ & 请求 $r$ 在 $t_n$ 是否已到站且仍在等待换电，由到站时间和此前换电次数确定 \\
$w_r$ & 请求 $r$ 在预测终点是否仍需服务，由第~\ref{sec:value}~节的式~\eqref{eq:terminal_pending} 和式~\eqref{eq:terminal_random} 确定 \\
$t_r^{\mathrm{arr}}$ & 请求的到站时刻；后站到站时刻随前站换电安排确定 \\
$e_{r,n},h_r$ & 分别表示 $t_n$ 是否处于到站至截止时间内，以及截止时刻是否早于预测终点；到站时间尚未确定时均为 $0$ \\
$I$ & 本轮预测期内的服务收入（元） \\
$C^{\mathrm{ch}},C^{\mathrm{adj}},C^{\mathrm{fail}}$ & 本轮预测期内的充电费用、本轮路径调整惩罚及本轮预测期内未满足预约请求的惩罚（元） \\
$s_\ell,s_{\ell+H}$ & 本轮观测运营状态与由本轮决策确定的预测终端状态 \\
$V_\theta(s)$ & 强化学习估计的状态终端价值（元） \\
\end{longtable}

\subsection{基于 SOC 分档的候选网络}
\label{sec:network}
对每个 O--D 对 $p$，以入口、沿线换电站和出口构造网络。车辆路径
$o^p\to i_1\to\cdots\to i_q\to d^p$
表示依次在 $i_1,\ldots,i_q$ 换电，弧所跨越的其他站点不被访问。将入口 SOC 划分为若干区间，分别预先生成候选网络，以减少在线需要处理的路径组合。

节点集合为 $\mathcal V^p=\{o^p\}\cup\mathcal S^p\cup\{d^p\}$，所有节点沿下游方向排列。以 $D_{i,j}^p$ 和 $v_{i,j}^p$ 表示节点间距离和行驶 SOC 消耗，以 $\underline s_d^p$ 表示出口 SOC 下限。最小换电间距 $\underline D$ 用于优先排除距离过近的首段和站间换电安排，出口末段按续航条件处理。对每个 SOC 区间 $h$，按以下四步生成候选网络。
\begin{enumerate}[label=\textbf{步骤\arabic*.},leftmargin=4.2em]
\item \textbf{生成原始弧。}若以该 SOC 区间上限电量能够从入口到达下游站点 $j$，则连接入口与 $j$；若 $v_{i,j}^p\le1$，则连接下游站点 $i,j$；若 $v_{i,d^p}^p+\underline s_d^p\le1$，则连接站点与出口。以该 SOC 区间上限电量从入口直接驶向出口后仍满足出口 SOC 下限时，加入相应弧。
\item \textbf{生成完整路径。}在原始网络中枚举所有入口至出口的完整路径。
\item \textbf{剪枝过近换电弧。}找出距离小于 $\underline D$ 的首段弧和站间弧，按距离从小到大考察。若当前仍存在至少一条不包含该弧的完整路径，则删除该弧；否则保留。
\item \textbf{形成候选网络。}合并最终保留路径中的弧，得到 SOC 区间对应的弧集 $\mathcal A^{h,p}$。
\end{enumerate}

在线使用时，根据用户实际入口 SOC 进一步检查首段续航。对在途用户，删除已经经过的节点，以实时位置建立虚拟起点，按实时 SOC 检查其到下游站点的首段弧。最小换电间距始终以入口或最近一次实际换电位置为基准。对已经在站等待的用户，将当前站作为后续路径起点，后续行程以该次换电完成后的满电状态计算；当前等待请求在第~\ref{sec:model}~节中保留，并与后续服务关联。

从本轮起点 $o^{p,k}$ 直接驶向出口 $d^p$ 的弧，只有满足以下相应条件时才予以保留：
\begin{equation}
\begin{cases}
s_{\mathrm{dep}}^{p,k}-v_{o^{p,k},d^p}^{p,k}\ge\underline s_d^p,
&\text{用户正在行驶或尚未进入高速},\\
1-v_{o^{p,k},d^p}^{p}\ge\underline s_d^p,
&\text{用户当前在站等待换电}.
\end{cases}
\label{eq:direct_exit_soc}
\end{equation}
式~\eqref{eq:direct_exit_soc} 第一行使用本轮起点的已知车辆 SOC $s_{\mathrm{dep}}^{p,k}$，其中 $v_{o^{p,k},d^p}^{p,k}$ 为从该起点直接驶向出口的 SOC 消耗。第二行使用当前站换电后的满电状态。两种情况都要求车辆到达出口时的 SOC 不低于 $\underline s_d^p$。由此得到用户弧集 $\mathcal A^{p,k}$、起点 $o^{p,k}$ 以及对应出口 $d^p$。

\subsection{日前换电路径计算}
\label{sec:dayahead}
日前换电路径根据预约用户的 O--D、入口 SOC 和候选网络计算。先使用实际预约入口 SOC 检查首段弧，并保留满足出口 SOC 要求的完整路径。设所得完整路径集合为 $\mathcal Y^{p,k}$，其元素由弧指示量 $y$ 表示，则日前路径取为
\begin{equation}
\bar y^{p,k}\in\arg\min_{y\in\mathcal Y^{p,k}}
\sum_{(j,i)\in\mathcal A^{p,k}:\,i\in\mathcal S}y_{j,i}.
\label{eq:dayahead_path}
\end{equation}
该目标选择换电次数最少的完整路径。当多条路径的换电次数相同时，依次比较其第一、第二及后续换电站的位置，优先选择更靠近出口的站点；站点位置相同的记录使用固定网络顺序区分。得到的 $\bar y^{p,k}$ 及对应换电站序列作为在线路径调整的基准。该计算仅需要车辆与道路信息，输入用户均具有满足上述条件的完整可达路径。

\subsection{滚动时域优化模型}
\label{sec:model}

\paragraph{请求与输入信息。}
本轮模型中的换电请求包括预约用户当前在站等待的请求、后续行程中可能产生的请求，以及随机请求。以 $r$ 标识一个换电请求，其所在站点记为 $i(r)$。令 $\R_i^A,\R_i^R$ 分别为站点 $i$ 的预约请求集合和随机请求集合，$\R_i=\R_i^A\cup\R_i^R$，并以 $\R^A,\R^R,\R$ 表示各站相应集合的并集。请求集合的具体构造、到站时间输入和退回 SOC 的计算方法在下文对应约束中说明。

\paragraph{决策变量与计算量。}
本轮优化选择预约用户的换电路径、各请求的换电时刻和所用电池，以及各充电槽的充电功率。用二元变量 $y_{i,j}^{p,k}\in\{0,1\}$ 表示用户 $(p,k)$ 的换电路径是否经过弧 $(i,j)$。一组满足流平衡的 $y$ 变量唯一确定该用户依次访问的换电站。用非负连续变量 $P_{i,b,n}\ge0$ 表示站点 $i$ 的充电槽 $b$ 在时间段 $n$ 内的充电功率。

以 $\alpha_{b,r,n}\in\{0,1\}$ 表示是否在 $t_n$ 用站点 $i(r)$ 的充电槽 $b$ 中的电池服务请求 $r$，其中 $b\in\mathcal B_{i(r)}$。电池在 $t_n$ 当次换电前的 SOC 记为 $S_{i,b,n}^{-}$，由初始电量、换电安排和充电功率递推得到。

下文用 $\chi_{p,k}$ 表示路径是否改变，$f_r$ 表示预约请求是否未获满足，$Q_{r,n}$ 表示请求是否正在等待换电，$w_r$ 表示请求在预测终点是否仍需服务。这些量均由相应定义确定，取值为 $0$ 或 $1$。$\chi_{p,k}$ 比较本轮选择的换电站与已有路径：未进入高速的用户与日前路径比较，在途用户与最近发布的路径比较。

$f_r=1$ 表示请求 $r$ 需要服务、等待截止时刻早于本轮预测终点，且运营商未能在等待期限内为其提供换电服务，即该预约请求未获满足。$Q_{r,n}=1$ 表示请求在 $t_n$ 已到站、尚未超过等待截止时刻，并且此前尚未换电。$w_r=1$ 表示预测终点仍需继续处理该请求。若某次预约请求超过等待期限仍未获服务，该用户的后续请求不再保留。

本轮开始时的日前路径、保留计划和最近发布路径均为已知输入。路径、换电和充电所需的计算量在对应约束中定义，终端状态相关量在第~\ref{sec:value}~节中定义；求解时所需的辅助变量和线性化按附录~\ref{lin:appendix} 处理。

\paragraph{目标函数。}

\begin{equation}
\max\quad I-C^{\mathrm{ch}}-C^{\mathrm{adj}}-C^{\mathrm{fail}}
+\beta V_\theta(s_{\ell+H}),
\label{eq:objective}
\end{equation}
其中，各项分别为
\[
\begin{aligned}
I&=E_B\sum_{n\in\mathcal T_\ell}\sum_{r\in\R}
c_{i(r),n}^{\mathrm{sw}}(1-\rho_r)\sum_{b\in\mathcal B_{i(r)}}\alpha_{b,r,n},\\
C^{\mathrm{ch}}&=\delta\sum_{n\in\mathcal T_\ell}\sum_{i\in\mathcal S}
c_{i,n}^{\mathrm{ch}}\sum_{b\in\mathcal B_i}P_{i,b,n},\\
C^{\mathrm{adj}}&=c^{\mathrm{adj}}\sum_{p\in\mathcal P}\sum_{k\in\mathcal K^p}\chi_{p,k},\\
C^{\mathrm{fail}}&=c^{\mathrm{fail}}\sum_{r\in\R^A}f_r.
\end{aligned}
\]

$I$ 为本轮预测期内的换电服务收入。每次换电交付满电电池，并收回 SOC 为 $\rho_r$ 的电池，因此按补充电量 $E_B(1-\rho_r)$ 与换电时刻对应的服务单价计费。

$C^{\mathrm{ch}}$ 为本轮预测期内的充电费用。其中，$c_{i,n}^{\mathrm{ch}}$ 表示站点 $i$ 在时间段 $n$ 的充电电价，充电槽 $b$ 在时间段 $n$ 内从电网输入的电量为 $\delta P_{i,b,n}$。

$C^{\mathrm{adj}}$ 为本轮优化中的路径调整惩罚。尚未进入高速公路的用户与日前路径比较，在途用户与最近一次向其发布的路径比较。只要新增或取消任意一个换电站，就对该用户计一次 $c^{\mathrm{adj}}$，改变多个站点也只计一次。

未进入高速的用户偏离日前路径时，该惩罚只用于本轮优化，不计入实际运营费用，也不在多轮之间累计结算。实际调整费用只在向在途用户发布不同于原安排的路径时计入，具体见式~\eqref{eq:actual_reward}。

$C^{\mathrm{fail}}$ 为本轮预测期内未满足预约请求的惩罚。当运营商未能在等待期限内满足预约换电请求时，计入一次惩罚 $c^{\mathrm{fail}}$。

$V_\theta(s_{\ell+H})$ 估计从预测终点状态出发的后续累计净收益。其输入 $s_{\ell+H}$ 和价值函数分别由第~\ref{sec:value}~节的式~\eqref{eq:terminal_state} 和式~\eqref{eq:value_network} 给出。它考虑终点时的电池电量和未完成换电请求，使本轮决策也考虑预测范围之外的充电与服务需求。$\beta$ 调整后续净收益在目标函数中的权重。


\paragraph{（1）流平衡约束。}
对每个用户，路径满足
\begin{equation}
\sum_{j:(i,j)\in\mathcal A^{p,k}}y_{i,j}^{p,k}
-\sum_{j:(j,i)\in\mathcal A^{p,k}}y_{j,i}^{p,k}
=\begin{cases}
1,&i=o^{p,k},\\-1,&i=d^p,\\0,&\text{其他节点},
\end{cases}
\quad i\in\mathcal V^{p,k},\ p\in\mathcal P,\ k\in\mathcal K^p.
\label{eq:flow}
\end{equation}
式~\eqref{eq:flow} 确定一条起点至出口的路径。

\paragraph{（2）路径调整约束。}
用 $x_i^{p,k}$ 表示本轮路径是否安排预约用户 $(p,k)$ 在站点 $i$ 换电：
\begin{equation}
x_i^{p,k}=\sum_{j:(j,i)\in\mathcal A^{p,k}}y_{j,i}^{p,k},
\qquad i\in\mathcal S,\quad p\in\mathcal P,\quad k\in\mathcal K^p.
\label{eq:station_visit}
\end{equation}


$\bar x_i^{p,k}$ 表示日前路径是否安排该用户在该站换电。对于尚未进入高速公路的用户，以 $\widehat x_i^{p,k}$ 表示本轮开始时保留的计划是否安排其在该站换电；首次参与优化时取 $\widehat x_i^{p,k}=\bar x_i^{p,k}$，以后使用上一轮保留的计划。对于在途用户，$x_i^{p,k,\mathrm{pub}}$ 表示最近一次向其发布的路径是否安排其在该站换电。已经完成的换电和当前正在等待的换电不参与路径调整，比较各条路径时同时去除这些站点；当前等待请求单独保留。各路径未安排换电的站点取 $0$。

两类用户都只能在式~\eqref{eq:update_clock} 规定的时刻调整路径，因此对各 $i\in\mathcal S$、$p\in\mathcal P$，有
\begin{subequations}\label{eq:path_freeze}
\begin{equation}
|x_i^{p,k}-\widehat x_i^{p,k}|\le g_\ell,
\qquad k\in\mathcal K^{\mathrm{fut},p}.
\label{eq:path_freeze_future}
\end{equation}
\begin{equation}
|x_i^{p,k}-x_i^{p,k,\mathrm{pub}}|\le g_\ell,
\qquad k\in\mathcal K^{\mathrm{enr},p}.
\label{eq:path_freeze_enroute}
\end{equation}
\end{subequations}
在 $g_\ell=0$ 时，式~\eqref{eq:path_freeze_future} 要求未进入高速的用户保持已保留的计划，式~\eqref{eq:path_freeze_enroute} 要求在途用户保持最近发布的路径。$g_\ell=1$ 时，两类用户都可以重新选择换电站。

路径调整惩罚按以下关系计算：
\begin{equation}
\chi_{p,k}=\begin{cases}
\displaystyle\max_{i\in\mathcal S}|x_i^{p,k}-\bar x_i^{p,k}|,
&k\in\mathcal K^{\mathrm{fut},p},\\
\displaystyle\max_{i\in\mathcal S}|x_i^{p,k}-x_i^{p,k,\mathrm{pub}}|,
&k\in\mathcal K^{\mathrm{enr},p},
\end{cases}
\qquad p\in\mathcal P.
\label{eq:path_change}
\end{equation}
式~\eqref{eq:path_change} 中，任意一个站点的换电安排与比较路径不同，就有 $\chi_{p,k}=1$。未进入高速的用户可以根据供需变化提前调整计划，但在进入高速前不发布，以避免反复通知用户。惩罚始终相对日前计划计算，以固定的比较依据限制不必要的路径调整。式~\eqref{eq:path_freeze} 和式~\eqref{eq:path_change} 的线性化方法见附录~\ref{lin:path_adjustment_section}。

\paragraph{（3）换电次数约束。}
对于用户 $k\in\mathcal K^p$，对每条到达换电站 $i$ 的可用弧 $(j,i)\in\mathcal A^{p,k}$，建立请求 $r=(p,k,j,i)$，其站点为 $i(r)=i$。本轮可用弧对应的预约请求组成集合
\[
\begin{aligned}
\R^{\mathrm{path}}=\bigl\{(p,k,j,i)\in\R:\;&
p\in\mathcal P,\ k\in\mathcal K^p,\\
&(j,i)\in\mathcal A^{p,k},\ i\in\mathcal S\bigr\}.
\end{aligned}
\]
$\R^{\mathrm{path}}$ 包含所有预约用户沿本轮可用弧可能产生的换电请求。对于已经在站等待的预约用户，当前站是本轮路径的起点，驶入该站的弧不在本轮可用弧集合中。该用户当前等待的请求单独保留，后续沿可用弧产生的请求仍属于 $\R^{\mathrm{path}}$。因此，$\R\setminus\R^{\mathrm{path}}$ 由当前在站等待的预约请求和全部随机请求组成。这些集合在本轮求解前确定。

用 $a_r$ 表示是否需要为请求 $r$ 安排换电，其取值为
\begin{equation}
a_r=\begin{cases}
y_{j,i}^{p,k},&r=(p,k,j,i)\in\R^{\mathrm{path}}\\
1,&r\in\R\setminus\R^{\mathrm{path}}
\end{cases}
\label{eq:request_selection}
\end{equation}
式~\eqref{eq:request_selection} 第一行规定：所选路径经过弧 $(j,i)$ 时，$a_r=1$；否则 $a_r=0$。第二行将当前在站等待的请求和随机请求的 $a_r$ 固定为 $1$。因此，$a_r$ 的取值由上述等式确定；取 $1$ 时能否完成换电，还需满足后续服务约束。

每个请求在本轮预测期内最多完成一次换电，同一充电槽在同一时刻最多交付一块电池。$\alpha_{b,r,n}=1$ 表示在 $t_n$ 用充电槽 $b$ 中的电池为请求 $r$ 换电。相应约束为
\begin{subequations}\label{eq:assignment_capacity}
\begin{align}
\sum_{n\in\mathcal T_\ell}
\sum_{b\in\mathcal B_{i(r)}}\alpha_{b,r,n}
&\le a_r\quad(r\in\R),\label{eq:request_service_limit}\\
\sum_{r\in\R_i}\alpha_{b,r,n}&\le1
\quad(i\in\mathcal S,b\in\mathcal B_i,n\in\mathcal T_\ell).
\label{eq:battery_service_limit}
\end{align}
\end{subequations}
式~\eqref{eq:request_service_limit} 左侧为请求 $r$ 在本轮预测期内的换电总次数，不包含预测终点 $t_{\ell+H}$ 当次的换电。$a_r=0$ 时，该请求不能获得服务；$a_r=1$ 时，最多获得一次服务。式~\eqref{eq:battery_service_limit} 将同一充电槽在 $t_n$ 服务的请求数限制为最多一个。其中，$\R$ 为本轮全部请求集合，$\R_i$ 为站点 $i$ 的请求集合，$\mathcal B_i$ 为该站的充电槽集合；约束对所有站点 $i\in\mathcal S$ 和本轮各时刻 $t_n$（$n\in\mathcal T_\ell$）成立。

\paragraph{（4）到站时刻与等待期限。}
用户到站后才能换电，超过等待截止时刻后不再安排服务。

以 $\R_{p,k}^A$ 表示预约用户 $(p,k)$ 在本轮模型中的所有换电请求，包括当前在站等待的请求，以及后续行程中可能产生的请求。对于请求 $r=(p,k,j,i)$，以 $\mathcal P_r$ 表示该用户在前一站 $j$ 的所有可能换电请求，即
\[
\mathcal P_r=
\begin{cases}
\bigl\{q\in\R_{p,k}^A:i(q)=j\bigr\}
&r\in\R^{\mathrm{path}}\\
\varnothing,&r\in\R\setminus\R^{\mathrm{path}},
\end{cases}
\]
其中，$q$ 表示同一用户的一个换电请求，$i(q)$ 为该请求所在的站点。用户可能通过不同路径到达前一站 $j$，所以 $\mathcal P_r$ 中可能包含多个请求，但所选路径至多采用其中一个。若用户当前已在 $j$ 等待换电，则在前往下一换电站之前，需要先完成当前这次换电。因此，下一站请求的 $\mathcal P_r$ 只包含当前在 $j$ 等待的请求。

当前等待请求自身和随机请求的 $\mathcal P_r$ 为空集。对于从入口或当前行驶位置出发的首次换电请求，本轮起点没有尚需服务的换电请求，因此 $\mathcal P_r$ 也为空集。$\mathcal P_r=\varnothing$ 表示在本轮安排请求 $r$ 换电之前，不需要先完成该用户的另一次换电。本轮开始前已经完成的换电不再包含在请求集合中。上述集合均在求解前确定。

对于 $\mathcal P_r=\varnothing$ 的请求，以 $\bar t_r^{\mathrm{arr}}$ 表示本轮给定的到站时刻。当前在站等待的请求使用实际到站时刻；用户当前未在站等待时的下一次换电请求，以及尚未到站的随机请求，使用本轮观测和预测给出的预计到站时刻。这些时刻在本轮求解中保持不变。下一轮更新时，已经到站的请求保留原始到站时刻，尚未到站的请求根据新观测更新预测。

对于 $\mathcal P_r\ne\varnothing$ 的请求，下一站的到站时刻需要由前一站的换电安排确定。以 $\tau_{j,i}^{p,k}>0$ 表示用户 $(p,k)$ 从换电站 $j$ 到下一换电站 $i$ 的预测行驶时间。该参数在本轮求解前给定，求解过程中保持不变，下一轮根据新观测更新。前一站等待时间的变化通过换电时刻传递到后一站。用 $\Delta$ 表示自到站起允许的最大等待时间，请求的到站时刻满足：
\begin{equation}
t_r^{\mathrm{arr}}=\begin{cases}
\bar t_r^{\mathrm{arr}},&\mathcal P_r=\varnothing,\\
\displaystyle\sum_{q\in\mathcal P_r}\sum_{m\in\mathcal T_\ell}
(t_m+\tau_{j,i}^{p,k})\sum_{b\in\mathcal B_j}\alpha_{b,q,m},
&\mathcal P_r\ne\varnothing.
\end{cases}
\quad r\in\R
\label{eq:arrival_eta}
\end{equation}
所选路径至多采用前一站的一个请求，且该请求最多换电一次。若前一站在 $t_m$ 换电，式~\eqref{eq:arrival_eta} 得到 $t_m+\tau_{j,i}^{p,k}$。到站时刻确定后，等待截止时刻为该时刻加上允许等待时间 $\Delta$。

用 $e_{r,n}$ 表示 $t_n$ 是否位于请求到站之后、等待截止之前。对所有 $r\in\R$ 和未来时刻 $n\in\{\ell+1,\ldots,\ell+H\}$，定义
\begin{equation}
e_{r,n}=\begin{cases}
\ind\{\bar t_r^{\mathrm{arr}}\le t_n\le\bar t_r^{\mathrm{arr}}+\Delta\},
&\mathcal P_r=\varnothing,\\
\displaystyle\sum_{q\in\mathcal P_r}
\sum_{\substack{m\in\mathcal T_\ell\\
t_m+\tau_{j,i}^{p,k}\le t_n\le t_m+\tau_{j,i}^{p,k}+\Delta}}
\sum_{b\in\mathcal B_j}\alpha_{b,q,m},
&\mathcal P_r\ne\varnothing.
\end{cases}
\label{eq:service_window}
\end{equation}
式~\eqref{eq:service_window} 第一种情况直接使用给定到站时刻。第二种情况只计入能够使用户不晚于 $t_n$ 到站、且到 $t_n$ 尚未超过截止的前站换电安排。前站没有换电时，求和为 $0$，后站不能提前获得服务。$e_{r,n}=1$ 只说明时间条件满足，能否换电还要看其他约束。式中的时间条件均由给定时刻和预测行驶时间确定，取值随前站服务变量变化。

当前的 $e_{r,\ell}$ 由真实到站和等待情况确定：只有已经在站等待且尚未超过等待截止时刻的请求取 $1$。在两个决策时刻之间到站的用户，从下一时刻参与服务安排，等待期限仍从真实到站时刻计算。未来时刻按式~\eqref{eq:service_window} 计算，预测终点只计算状态。

\paragraph{（5）等待和服务顺序。}
本组确定每个时刻哪些请求正在等待换电，并规定多人同时等待时的服务优先顺序。等待标记和服务安排满足
\begin{subequations}\label{eq:queue}
\begin{equation}
Q_{r,n}=a_re_{r,n}\left[
1-\sum_{\substack{m\in\mathcal T_\ell\\m<n}}
\sum_{b\in\mathcal B_{i(r)}}\alpha_{b,r,m}
\right],
\qquad r\in\R,\quad n\in\mathcal T_\ell\cup\{\ell+H\}.
\label{eq:queue_status}
\end{equation}
\begin{equation}
\sum_{b\in\mathcal B_{i(r)}}\alpha_{b,r,n}\le Q_{r,n},
\qquad r\in\R,\quad n\in\mathcal T_\ell.
\label{eq:queue_service}
\end{equation}
\end{subequations}
式~\eqref{eq:queue_status} 中，求和项只计算从 $t_\ell$ 到 $t_n$ 之前已完成的换电次数，不包含 $t_n$ 当次的换电；$n=\ell$ 时求和为 $0$。请求需要服务、已经到站且尚未超过截止时刻，并且此前尚未换电时，才取 $Q_{r,n}=1$。因此，在 $t_n$ 获得服务的请求，本次仍有 $Q_{r,n}=1$。式~\eqref{eq:queue_service} 只允许为正在等待的请求换电，但不要求每个等待请求都立即获得服务。预测终点只计算 $Q_{r,\ell+H}$，不安排当次换电。

对于需要先在前一站换电的请求，式~\eqref{eq:arrival_eta} 与式~\eqref{eq:service_window} 已保证：只有前一站换电并经过正的行驶时间后，$e_{r,n}$ 才能取 $1$。因此，前站尚未换电时不能服务后站，同一用户也不能在同一时刻完成前后两次换电。式~\eqref{eq:queue_status} 的乘积关系按附录~\ref{lin:waiting_priority_section} 线性化。

在同一站点定义 $q\prec r$：当 $q$ 为预约而 $r$ 为随机请求，或者二者同类且 $t_q^{\mathrm{arr}}<t_r^{\mathrm{arr}}$ 时成立。同类请求的先后顺序由式~\eqref{eq:arrival_eta} 中本轮确定的到站时刻比较得到。对任意两个不同请求，施加以下条件约束：
\begin{equation}
\begin{gathered}
q\prec r\quad\Longrightarrow\quad
\sum_{b\in\mathcal B_i}\alpha_{b,r,n}
\le 1-Q_{q,n}+\sum_{b\in\mathcal B_i}\alpha_{b,q,n},\\
q,r\in\R_i,\quad q\ne r,\quad i\in\mathcal S,\quad n\in\mathcal T_\ell.
\end{gathered}
\label{eq:priority}
\end{equation}
式~\eqref{eq:priority} 要求：较高优先级请求仍在等待时，只有它在同一时刻获得服务，才能服务较低优先级请求。尚未到站、已经服务或已经失效的请求，其 $Q_{q,n}=0$，不会阻止其他请求服务。同类请求同时到站时不规定二者的先后。每个到站时刻只可能取给定时刻、有限个 $t_m+\tau_{j,i}^{p,k}$ 或未确定时的零值，因此上述严格先后关系可按有限个已知时刻的组合精确展开。到站时间尚未确定的请求有 $Q_{q,n}=0$，其填零的时刻不会使它优先占用服务。运营者可以在等待期限内选择换电时刻，并通过 $\alpha_{b,r,n}$ 确定各请求使用的电池。式~\eqref{eq:priority} 的等价线性约束见附录~\ref{lin:waiting_priority_section}。

\paragraph{（6）未满足预约请求的判定。}
用 $h_r$ 表示请求的等待截止时刻是否早于本轮预测终点，直接根据给定到站时间或前站换电安排计算：
\begin{equation}
h_r=\begin{cases}
\ind\{\bar t_r^{\mathrm{arr}}+\Delta<t_{\ell+H}\},
&\mathcal P_r=\varnothing,\\
\displaystyle\sum_{q\in\mathcal P_r}
\sum_{\substack{m\in\mathcal T_\ell\\
t_m+\tau_{j,i}^{p,k}+\Delta<t_{\ell+H}}}
\sum_{b\in\mathcal B_j}\alpha_{b,q,m},
&\mathcal P_r\ne\varnothing,
\end{cases}
\qquad r\in\R.
\label{eq:deadline_passed}
\end{equation}


用 $f_r$ 表示预约请求是否未获满足，其取值由
\begin{equation}
f_r=a_rh_r\left(
1-\sum_{n\in\mathcal T_\ell}
\sum_{b\in\mathcal B_{i(r)}}\alpha_{b,r,n}
\right),\qquad r\in\R^A
\label{eq:failure_exact}
\end{equation}
确定。式~\eqref{eq:failure_exact} 表示：只有请求需要服务、截止时刻早于预测终点，且仍未完成换电时，才取 $f_r=1$。其中的求和项为该请求在本轮预测期内的换电总次数。已经按时完成换电的请求，该求和项为 $1$，因此即使 $h_r=1$，仍有 $f_r=0$。式~\eqref{eq:queue_status} 和式~\eqref{eq:queue_service} 将换电限制在 $e_{r,n}=1$ 的时刻，因此不允许在等待截止之后服务。本组变量乘积的线性化见附录~\ref{lin:unmet_requests_section}。

若某站的预约请求超过等待期限仍未获服务，后续站点便因前站未换电而无法确定到站时刻。由式~\eqref{eq:deadline_passed} 和式~\eqref{eq:failure_exact}，这些后续请求的 $h_r$ 和 $f_r$ 均为 $0$。此前已完成换电的请求和未选择的请求也有 $f_r=0$，因此每位用户在本轮预测期内至多有一个预约请求计为未满足，惩罚不会重复计算。

若某站的预约请求在 $t_n$ 之前已超过等待期限且未获服务，该请求因超过截止时刻而有 $e_{r,n}=0$。在该用户所选路径上，后续请求因前站未换电而有 $e_{r,n}=0$，更早的请求已经服务，式~\eqref{eq:queue_status} 中的累计换电次数为 $1$，因此 $Q_{r,n}=0$；未选择的请求则有 $a_r=0$。这些条件已使该用户的请求退出等待，无需再逐时刻定义该用户是否已有未满足的预约请求。

截止不早于预测终点时，不能因为本轮尚未安排换电就提前将该预约请求判为未满足。仍需服务的请求保留在终端状态，由后续滚动优化继续安排。

\paragraph{（7）满电交付和充电。}
在每个时刻 $t_n$，先换电，再在随后的时间段 $\mathcal I_n$ 内充电。换电时收回的电池 SOC 记为 $\rho_r$。对于由可用弧建立的预约请求 $r=(p,k,j,i)\in\R^{\mathrm{path}}$，按车辆出发时的电量和行驶 SOC 消耗计算：
\[
\rho_r=
\begin{cases}
s_{\mathrm{dep}}^{p,k}-v_{o^{p,k},i}^{p,k},&j=o^{p,k},\ \text{用户正在行驶或尚未进入高速},\\
1-v_{j,i}^{p},&j\text{ 为换电站}.
\end{cases}
\]
其中，$s_{\mathrm{dep}}^{p,k}$ 为车辆在本轮起点的已知 SOC，$v_{o^{p,k},i}^{p,k}$ 为从该起点到站点 $i$ 的已知 SOC 消耗。尚未进入高速的用户使用预约信息及入口前观测，在途用户使用实时值。从换电站 $j$ 出发时使用换电后的满电状态，因此退回 SOC 为 $1-v_{j,i}^{p}$。用户当前在站等待时，后续首段也按这一情况计算。

当前在站等待请求的 $\rho_r$ 根据观测给定，尚未到站的随机请求使用本轮给定的预测值。所有请求的 $\rho_r$ 均在求解前确定；可用弧对应请求的取值由第~\ref{sec:network}~节的网络生成规则保证，其余请求的输入同样满足 $0\le\rho_r<1$。

以 $S_{i,b,n}^{-}$ 表示充电槽 $b$ 中电池在 $t_n$ 当次换电前的 SOC，将换电与随后充电的影响直接写入下一时刻的电量：
\begin{subequations}\label{eq:battery_dynamics}
\begin{align}
S_{i,b,\ell}^{-}&=S_{i,b}^{\mathrm{obs}},\label{eq:initial_soc}\\
S_{i,b,n}^{-}&\ge\alpha_{b,r,n},
&&r\in\R_i,\label{eq:full_delivery}\\
S_{i,b,n+1}^{-}&=S_{i,b,n}^{-}
-\sum_{r\in\R_i}(1-\rho_r)\alpha_{b,r,n}
+\frac{\eta_i\delta}{E_B}P_{i,b,n},\label{eq:charge_soc}\\
0\le S_{i,b,n}^{-}&\le1,
&&n\in\mathcal T_\ell\cup\{\ell+H\},\label{eq:soc_bounds}\\
0\le P_{i,b,n}&\le\overline P_{i,b}^{\mathrm{ch}},\label{eq:device_power}\\
\sum_{b\in\mathcal B_i}P_{i,b,n}&\le\overline P_i^{\mathrm{sta}}.\label{eq:station_power}
\end{align}
\end{subequations}
上述关系对所有 $i\in\mathcal S$ 和 $b\in\mathcal B_i$ 成立；站点总功率约束对各站点成立。除初始时刻和式~\eqref{eq:soc_bounds} 另有规定外，时间索引均取 $n\in\mathcal T_\ell$。

式~\eqref{eq:initial_soc} 将初始 SOC 固定为当前观测值。式~\eqref{eq:full_delivery} 保证满电交付：$\alpha_{b,r,n}=1$ 时，结合式~\eqref{eq:soc_bounds} 得到 $S_{i,b,n}^{-}=1$；$\alpha_{b,r,n}=0$ 时，不要求该电池满电，也不要求满电电池必须交付。

式~\eqref{eq:charge_soc} 中，下一时刻的 SOC 等于当前 SOC，减去换电造成的下降量，再加上本时段充电带来的增量。由式~\eqref{eq:battery_service_limit}，同一充电槽在同一时刻最多服务一个请求。若服务请求 $r$，满电电池被取走，退回电池的 SOC 为 $\rho_r$，因此扣减 $1-\rho_r$；若没有换电，求和项为 $0$。由于 $0\le\rho_r<1$，换电后槽内电池的 SOC 自然位于 $0$ 与 $1$ 之间。

充电功率 $P_{i,b,n}$ 是从电网输入的功率，$\delta P_{i,b,n}$ 为输入电量，乘以充电效率 $\eta_i$ 后得到存入电池的电量，再除以容量 $E_B$，得到 SOC 增量。功率以 kW 计、容量以 kWh 计，因此 $\delta$ 以小时计。例如，容量为 $60$~kWh、效率为 $0.9$、功率为 $40$~kW 时，充电 $5$~min 使 SOC 增加 $0.05$，即 $5$ 个百分点；若换电后电池的 SOC 为 $0.30$，下一时刻就是 $0.35$。

式~\eqref{eq:soc_bounds} 覆盖本轮各换电时刻和预测终点，保证电量非负且不超过容量。终点的 SOC 已包含最后一个时间段的充电结果，因此同样需要这一上界。

式~\eqref{eq:device_power} 限制单个充电槽的功率，取 $0$ 表示该时间段不充电。所选功率在该时间段内保持恒定，并与 SOC 上界共同限制充电量。式~\eqref{eq:station_power} 限制同一站点各充电槽的总功率。例如，两台设备的功率上限各为 $60$~kW，而站点上限为 $80$~kW 时，两台设备不能同时以 $60$~kW 充电。每个请求可以选择任意满足条件的充电槽，模型同时确定其后续充电功率。

式~\eqref{eq:flow}--\eqref{eq:battery_dynamics}、变量范围，以及第~\ref{sec:value}~节定义的终端状态和价值函数共同构成优化问题。到站、等待和截止条件由给定时刻及离散服务安排精确确定；服务先后关系按有限个候选到站时刻展开。再将绝对值、最大值、价值计算及有界变量乘积按附录~\ref{lin:appendix} 转换为线性约束，得到混合整数线性优化模型。

\paragraph{滚动执行与记账。}
在 $t_\ell$ 读取真实在站请求、电池状态、车辆位置、保留的计划和最新预测，建立本轮模型。路径更新时刻同时优化两类预约用户的路径。未进入高速的用户只保存优化后的计划，用于后续预测和下一轮的 $\widehat x_i^{p,k}$，不发布也不计入实际费用。在途用户的路径若与最近发布的安排不同，则发布新路径、计入一次实际调整费用，并将新路径作为后续比较的依据。非路径更新时刻，两类用户都保持各自已有的安排。

按模型确定的请求和电池安排执行当前换电后，在 $\mathcal I_\ell$ 内使用当前充电功率。区间内到站请求在下一时刻参与换电安排，原到站时刻用于计算等待时间。截止位于 $\mathcal I_\ell$ 内且尚未完成换电的请求在该时段判定超时；恰在下一时刻截止的请求留到下一轮先参与换电安排。

下一轮电池状态由实际服务和充电更新，真实等待请求保留原标识及截止时间。车辆位置、下一次换电的预计到站时刻和站间行驶时间根据新观测更新；仍需先完成前站换电的后续请求继续由式~\eqref{eq:arrival_eta} 计算到站时刻。预测请求在实际到站后使用真实到站时刻，状态中始终区分已观测请求与预测请求。若实际执行中，前一站的预约请求超过等待期限仍未获服务，该用户的后续服务也随之取消。新的当前状态成为下一轮初始条件，新的路径更新参数按式~\eqref{eq:update_clock} 计算。
